\documentclass[12pt]{article}
\usepackage[utf8]{inputenc}
\usepackage[T1]{fontenc}
\usepackage{hyperref}
\title{Kven Finnish Language Resources for K-12 Curriculum Development}
\author{Compiled by Claude Sonnet 5.5 (claude-sonnet-5-5) for Andr\'as Kornai}
\date{}
\begin{document}
\maketitle

\section{Introduction}
Kven Finnish (ISO 639-3 fkv; Glottocode kven1236) is a Uralic (Finnic) language of northern Norway, recognised in Norway as a language in its own right in 2005. This survey is based on a small number of web searches; it lists only resources that appeared in search results.

\section{Primary Educational Resources}
No Kven textbook or school primer was found in the search results. Classroom materials may exist (for example from the Kven Institute) but were not located and are not claimed here.

\section{Linguistic and Reference Materials}
\begin{itemize}
\item Norwegian Bokm\r{a}l--Kven dictionary (Giellatekno, UiT, with Kainun institutti; 7,716 entries, open licence): \url{https://b2find.eudat.eu/dataset/595deefa-81cf-530c-b5dc-ace54148f680}
\item Kvensk revitalisering, normering og leksikografi (academic work, UiT Munin): \url{https://munin.uit.no/handle/10037/36648}
\item Article on the online Kven--Norwegian dictionary (LexicoNordica): \url{https://tidsskrift.dk/lexn/article/download/151290/193992/333797}
\item Spr\r{a}kbanken resource catalogue entries (CLARINO deposits): \url{https://www.nb.no/sprakbanken/ressurskatalog/oai-repo-clarino-uib-no-11509-92/} and \url{https://www.nb.no/sprakbanken/ressurskatalog/oai-repo-clarino-uib-no-11509-89/}
\end{itemize}

\section{Religious and Community Resources}
\begin{itemize}
\item Norwegian Kven Association, ``Kven language and culture'': \url{https://kvener.no/?p=5939}
\end{itemize}

\section{Digital and Audio Resources}
The Kven Institute and Giellatekno developed an online dictionary and language technology; the sources report that several hundred hours of old spoken recordings have not yet been used. See the dictionary and LexicoNordica entries above.

\section{Curriculum Development Context}
Dictionary and language-technology infrastructure exists, but no classroom curriculum or textbook was found in these searches. Any K-12 curriculum would need to build on the dictionary, the normative work on a written standard, and the unused spoken archives.
\end{document}